% Sections 1-9: production, flows, boundary, residence, output, prices, external trade, OS, D&D.

%======================================================================
\section{Production, value added and the circular flow}\label{sec:prod}
\covers{T1, T2, T3}{Lectures 1--2}{2025--26 Q1(i), (ii)}

Production is measured without double counting. Marshall's example is a carpet: once the carpet
is counted at full value, the yarn and the labour in it are already counted. A producer is
therefore credited only with the value it \emph{adds}.

\begin{tabularx}{\textwidth}{@{}>{\bfseries}l X@{}}
IC & \df{Intermediate consumption}: goods and services a producer buys and \emph{uses up during} the
production process of the period, such as raw materials, fuel, power, maintenance and hired services. \\
GVO & \df{Gross value of output}: the value of all goods and services the unit produces in the period. \\
GVA & \df{Gross value added} $=$ GVO $-$ IC: the unit's own contribution to production. \\
GDP & \df{Gross domestic product}: the sum of GVA over all resident producers, plus net product taxes.
This is the value of production within the economic territory. \\
HFCE & \df{Household final consumption expenditure}: what resident households spend on goods and
services for final consumption. \\
GCF & \df{Gross capital formation}, or investment: acquisitions less disposals of produced assets,
namely fixed assets (GFCF), change in inventories and valuables. \\
GNI & \df{Gross national income} $=$ GDP $+$ net primary (earned) income from the rest of the world.
This is the income of residents. \\
\end{tabularx}

\begin{keybox}{Relations to reproduce}
\vspace{-8pt}
\begin{align*}
\T{GVA} &= \T{GVO}-\T{IC}, \qquad \T{GVA}_{bp}=\T{GVO}_{bp}-\T{IC}\\
\T{GDP} &= \textstyle\sum \T{GVA}_{bp} + (t-s)_{\text{products}} + (t-s)_{\text{imports}}\\
\T{GVO} &= \T{HFCE}+\T{IC}+\T{GCF}\ \Longrightarrow\ \T{GVA}=\T{HFCE}+\T{GCF},\ \text{i.e. } Y=C+I
  \quad(\text{closed, no govt.})\\
\T{GNI} &= \T{GDP} + \text{net primary income from RoW}, \qquad \text{Net} = \text{Gross} - \T{D\&D}
\end{align*}
\end{keybox}

\subsection*{The two-sector circular flow (T2)}

\begin{center}
\begin{tikzpicture}[font=\small,>=Stealth,
  box/.style={draw,thick,minimum width=2.6cm,minimum height=0.8cm}]
\node[box,fill=Yellow!35] (H) at (0,0) {Households};
\node[box,fill=Lavender!45] (E) at (9,0) {Enterprises};
\draw[->,very thick,Blue] (E.north) to[out=150,in=30]
  node[midway,above]{goods \& services} (H.north);
\draw[->,very thick,BrickRed] ([yshift=6pt]H.east) --
  node[midway,above]{expenditure on g\&s} ([yshift=6pt]E.west);
\draw[->,very thick,BrickRed] ([yshift=-6pt]E.west) --
  node[midway,below]{factor compensations (w\&s, OS)} ([yshift=-6pt]H.east);
\draw[->,very thick,Blue] (H.south) to[out=-30,in=-150]
  node[midway,below]{labour and other factor services} (E.south);
\node[font=\footnotesize,align=left,anchor=west] at (10.5,0)
  {\textcolor{Blue}{blue: real flows}\\ \textcolor{BrickRed}{red: money flows}};
\end{tikzpicture}
\end{center}

\begin{itemize}
\item Households own \emph{all} factors of production (land, labour, capital, entrepreneurship). They
  sell factor services and are paid factor compensations: wages \& salaries and operating surplus.
  Those payments are the economy's income.
\item Real flows (goods and services, factor services) run opposite to money flows (expenditure,
  compensations).
\item The model rests on Say and Walras: there is no lag between production and income, and no
  income goes unspent. So income equals expenditure (Say's law, ``supply creates its own demand'').
  Keynes (1936) rejected this: spending depends on income, effective demand drives output, and
  aggregate expenditure is $C+I$.
\item Hence \textbf{production $\equiv$ income $\equiv$ expenditure}, which gives the three ways of
  measuring GDP. The \emph{production} approach sums GVA over all producers. The \emph{income}
  approach sums the incomes generated. The \emph{expenditure} approach sums \emph{final}
  expenditure on the goods and services produced.
\end{itemize}

\subsection*{\texorpdfstring{$Y = C + I$}{Y = C + I} (T3)}
The simple circle leaves out enterprise-to-enterprise flows: inputs (IC) and capital goods
(investment). Once they are added, the inputs cancel in aggregation. What remains is
consumption goods plus capital goods, so $Y \equiv C + I_g$. Output left unsold goes into
inventories. The change in inventories (CII) is part of production \emph{and} of investment, so
the identity still holds, and so does $I_g \equiv S$.

\begin{tmpl}{Lecture 2 illustration}
\small Sales 18,000, made up of inputs 3,000, consumption goods 14,000 and capital goods
(tractors) 1,000. Payments: rent 1,200, interest 900, wages \& salaries 8,900.\\
Production $Y = 18{,}000 - 3{,}000 = 15{,}000 = C\,(14{,}000) + I\,(1{,}000)$. Income
$= 8{,}900 + 1{,}200 + 900 + \text{profit }4{,}000 = 15{,}000$, where profit is the residual.\\
\emph{Twist:} households spend only 13,500. Then $S = 15{,}000 - 13{,}500 = 1{,}500$, the 500
unsold goes into CII, and $I_g = 1{,}000 + 500 = 1{,}500 = S$.\\[2pt]
\textbf{Symbolic.} Sales $V$, inputs $N$, consumption goods $C_g$, capital goods $K$:
$Y = V-N = C_g+K$. If households spend $C<C_g$: $\T{CII}=C_g-C$, $S=Y-C$, and $I_g=K+\T{CII}=S$.
\end{tmpl}

\begin{quick}{which are true?}
\small (a) GDP can be measured by production, income and expenditure \yes; (b) production approach
$=$ GVA of all units \yes; (c) income approach $=$ sum of incomes of all households \yes;
(d) expenditure approach $=$ sum of \emph{all} expenditure on goods and services \no, because that
includes IC and only final expenditure counts; (e) total expenditure $=$ FCE $+$ GCF $+$ exports
net of imports \yes.
\end{quick}

%======================================================================
\section{Economic flows and transactions}\label{sec:flows}
\covers{T4, T5, T6, T25}{Lecture 3}{2025--26 Q1(viii)}

An \df{economic flow} is any change in the volume or value of a unit's economic assets. It can be
caused by human acts, by changes in prices, or by natural events. Examples are production, sale
and purchase, taxes, transfers, destruction by calamity, and holding gains.

\begin{multicols}{2}\small
\textbf{Economic flows}
\begin{itemize}
\item \textbf{Transactions}: payments and receipts for production, consumption and accumulation,
  and transfers.
  \begin{itemize}
  \item \df{Exchange}: something is given in return.
  \item \df{Transfer}: nothing is given in return.
  \end{itemize}
\columnbreak
\item \textbf{Other flows}
  \begin{itemize}
  \item \emph{Other changes in volume}: economic value appears or disappears (discoveries,
    destruction by calamity).
  \item \emph{Changes in the level and structure of prices}: holding gains and losses.
  \end{itemize}
\end{itemize}
\end{multicols}

A \df{transaction} is an economic flow in which institutional units interact by mutual agreement.
It can also be an action within one unit that is analytically useful to treat as a transaction,
usually because the unit is acting in two capacities. Transactions are either actual (observed)
or notional (not observable, so imputed). They cover production and exchange of products and
non-produced assets, distribution of income, transfers (current and capital), and financial
assets and liabilities.

\begin{keybox}{Current vs capital transactions (asked: 2025--26 Q1(viii))}
\small
\df{Current transactions} relate to production, final consumption and current transfers:
production-related flows (output, IC, CFC), primary income (CE, OS/MI), production taxes less
subsidies, final consumption, and current transfers such as income taxes and transfers in money
or kind. They do not change the stock of assets directly and are recorded in the \emph{current
accounts}.\\[2pt]
\df{Capital transactions} create or change the volume of a unit's assets: capital formation,
acquisition and disposal of non-produced (natural) assets, capital taxes and capital transfers.
They are recorded in the \emph{accumulation accounts}, and they change the balance sheet.
\end{keybox}

\df{Other perspectives} (T25). A transaction is \emph{between} units or \emph{within} one unit;
within-unit transactions are always non-monetary. It is \emph{monetary} if money changes hands,
and a monetary transaction always has two parties. \emph{Non-monetary} transactions have no
money flow, so they cannot be observed and are imputed, usually at market value. See
\S\ref{sec:mta} for the compiler's tree.

\begin{quick}{economic flows? transactions?}
\small \textbf{Economic flow?} Paddy cultivation \yes; destruction of property by a calamity \yes;
father giving money to his son in the same household \no\ (within one unit); drawing river water
to irrigate land \yes; natural growth of fish stock in a river \no\ (not an economic asset);
price appreciation of real estate \yes.\\
\textbf{Transaction: exchange or transfer?} Production of crops: exchange. Income tax: transfer.
Sale of smuggled goods: exchange (illegal trade still counts). Pickpocketing: transfer. Paying a
lost bet: transfer.
\end{quick}

%======================================================================
\section{The production boundary and IPP}\label{sec:boundary}
\covers{T7, T8}{Lecture 4}{---}

\df{Production} (SNA) is a physical process, under the control and responsibility of an
institutional unit, that uses labour and assets to turn inputs of goods and services into
outputs. The output must be capable of being supplied to another unit, with or without charge.
There are three kinds of product. \emph{Goods} are physical, in demand, and carry ownership that
can be transferred. \emph{Services} change the condition of the consumer or help exchange
products or financial assets (traders, banks). \emph{Data} (2025 SNA) is information produced by
observing phenomena and stored digitally; it is capitalised if used for more than a year, and
then depreciated.

\df{General production boundary}: every activity that one unit could pay another to do for it,
the ``third-person'' test. It excludes only basic human activities such as eating, drinking,
sleeping and taking exercise, which nobody else can do on your behalf.

\begin{keybox}{SNA (2008/2025) production boundary --- inclusion and exclusion}
\small
\begin{tabularx}{\linewidth}{@{}X|X@{}}
\textbf{Included} & \textbf{Excluded} \\ \hline
All \emph{goods}, whether for the market or for own final use (a farmer's crops eaten at home). &
Own-account \emph{household services} by unpaid members: cooking, cleaning, care, repair of own
durables, transporting members. \\
\emph{Services} only if they are sold, supplied to other units, or produced by paid labour (paid
domestic staff). & Volunteer, social and cultural activities. \\
Housing services of owner-occupiers (imputed rent); own-account IPP. & Things that only affect
the quality of life. \\
Non-market output of government and NPISHs; bank services not explicitly charged (FISIM). &
\emph{Theft} (a transfer, not production). \\
Illegal and concealed production (narcotics trade); natural growth of \emph{cultivated} forests. &
Natural growth that nobody manages (wild fish, wild forests). \\
\end{tabularx}
\end{keybox}

\df{Intellectual property products} (T8) are the results of research, innovation and artistic
work. They have been fixed assets since the 2008 SNA. There are four categories: (1) research
and development; (2) mineral exploration and evaluation; (3) computer software, data and
databases; (4) entertainment, literary and artistic originals.

\begin{quick}{inside the SNA boundary? is it production?}
\small \textbf{Inside?} Maintenance and repair of household durables by the household \no; crops
kept for own consumption \yes; owner-occupied housing services \yes; transporting household
members \no; natural growth of cultivated forests \yes; distribution of narcotics \yes.\\
\textbf{Production?} A morning cup of tea \no\ (nothing is created); ploughing \yes; newspaper
delivery \yes\ (trade); repairing a vending machine \yes\ (adds value); natural growth of fish in
a river \no\ (no human effort); collecting wild honey \yes\ (human effort).
\end{quick}

%======================================================================
\section{The domestic economy: territory, units, residence}\label{sec:res}
\covers{--- (background for GNI, exports and imports)}{Lecture 5}{2025--26 Q1(i)}

\begin{itemize}
\item The \df{economy} is all institutional units \emph{resident} in the economic territory during
  the period.
\item The \df{economic territory} is the area the government administers, within which persons,
  goods and capital circulate freely. It includes airspace, territorial waters and continental
  shelf; the country's enclaves abroad (embassies, consulates, military bases, scientific
  stations); and free zones and offshore activity.
\item An \df{institutional unit} can own assets, incur liabilities, carry out economic activity
  and transact with others in its own right. All transactions take place through institutional
  units. A household that also produces is treated as two units, one producing and one consuming.
\item There are \df{five sectors}: non-financial corporations, financial corporations, general
  government, households (including unincorporated enterprises and the informal sector), and
  NPISHs. Government enterprises run like businesses (railways, airlines, power) are
  quasi-corporations and belong in the corporate sector. The central bank is a financial
  corporation.
\item \df{Residence} is the territory where a unit has its \emph{centre of predominant economic
  interest}. It does not depend on nationality or currency. An individual takes the residence of
  their household, and an unincorporated enterprise that of its owner. A corporation or NPI is
  resident where it is registered, and a foreign branch is a resident quasi-corporation of the
  host country. Land and buildings owned by non-residents are assigned a notional resident unit.
  Non-residents that transact with residents form the \df{rest of the world (RoW)}.
\item \df{Exports and imports} are changes of economic ownership between a resident and a
  non-resident.
\end{itemize}

\begin{quick}{export, import or neither (for India)? true or false?}
\small A Kolkata resident buys a Swiss watch in a Kolkata shop: \textbf{N} (the shop is
resident). Online purchase of a UK book: \textbf{M}. An Indian magazine delivered to an Indian
student at Oxford: \textbf{N} (the student is still resident). An Egyptian's hotel bill in
Mumbai: \textbf{X}.\\
Foreign students staying three years are residents \no. A Citibank branch in Kolkata is resident
in India \yes. The Australian crew of a Japanese ship are residents of Japan \no. Non-residents
are not treated as owners of immovable assets \yes. Unincorporated businesses without separate
accounts belong to households \yes. A foreign NPI's branch serving residents is a resident NPISH
\yes. The central bank is in general government \no. A Japanese firm's branch in Thailand is
resident in Thailand \yes.
\end{quick}

%======================================================================
\section{Economically significant prices and the three kinds of output}\label{sec:output}
\covers{T9, T13, T14}{Lectures 5, 7}{---}

\df{Economically significant prices} are prices that significantly affect how much producers are
willing to supply and how much purchasers want to buy. In practice, sales cover the majority of
production costs. General government, the central bank and NPISHs mostly charge nothing or prices
that are not economically significant.

\df{Output} (2025 SNA) is the goods and services an establishment produces. It excludes goods and
services used in an activity where the establishment bears no risk, and those the establishment
consumes itself, unless they are used for own capital formation or own final consumption.

\begin{center}\small
\begin{tabularx}{\textwidth}{@{}l X X X@{}}
\toprule
& \textbf{Market output} & \textbf{Non-market output} & \textbf{Output for own final use} \\
\midrule
What & Sold at economically significant prices & Supplied free, or at prices that are not
economically significant, to other units or the community & Kept by the producer for its own
final consumption or capital formation \\
Who & Corporations, households & General government, NPISHs, central bank (2025 SNA) & Any
producer \\
Valued at & Quantity $\times$ market price & Sum of costs & Market-equivalent price, else sum
of costs \\
Examples & Rented housing, transport corporations, restaurants & Police, free schools run by
trusts, government hospitals & A farmer's own crops, a retailer's own shelves, owner-occupied
housing, paid domestic staff, CII kept for these uses \\
\bottomrule
\end{tabularx}
\end{center}

\df{Sum of costs} (2025 SNA) $=$ IC $+$ remuneration of employees $+$ depreciation (and
depletion) $+$ other taxes less subsidies on production $+$ net return to non-financial assets
$+$ rent on non-produced non-financial assets. Products end up in one of three uses: IC,
consumption of fixed capital (CFC), or \df{final use}, which is final consumption or capital
formation.

\begin{quick}{economically significant? which kind of output?}
\small \textbf{ESP?} Entry ticket to a government museum \no; a pen at the local store \yes; a
restaurant meal \yes; a token OPD registration fee at a free government hospital \no; compulsory
payments by financial corporations to the central bank \no.\\
\textbf{Kind?} Police: non-market. Shelves a retailer builds for his shop: own final use.
Owner-occupied housing: own final use. Rented housing: market. Free schools run by trusts:
non-market. Transport corporations: market.
\end{quick}

%======================================================================
\section{Basic, producers' and purchasers' prices}\label{sec:prices}
\covers{T11}{Lecture 6}{2025--26 Q1(iii)}

A user usually pays more than the producer receives, because of taxes, trade margins and
transport costs. The SNA therefore fixes three valuations:
\begin{itemize}
\item \df{Purchasers' price}: what the purchaser pays to take delivery at the time and place
  required, including trade and transport margins (TTM) and taxes.
\item \df{Producers' price}: what the producer receives from the purchaser, excluding separately
  invoiced VAT and margins.
\item \df{Basic price}: what the producer finally keeps, i.e.\ net of taxes on products and
  including subsidies on products.
\end{itemize}

\begin{keybox}{The price chain}
\centering
Purchasers' price $-$ TTM (and taxes on consumers) $=$ Producers' price\\
Producers' price $-$ (taxes $-$ subsidies) on products $=$ Basic price\\[2pt]
\raggedright\small Output is recorded at \textbf{basic} prices; uses (IC and final use) at
\textbf{purchasers'} prices. The 2008 SNA avoids the term ``market prices'' for value added,
because VAT makes it ambiguous. Write $\T{GDP}=\T{GVA}_{bp}+(t-s)_{\text{products}}$ rather than
``GDP at market prices''.
\end{keybox}

\begin{tmpl}{Lecture 6 pen}
\small A consumer pays 20 to a trader. The trader bought the pen at 12 from the manufacturer,
who paid a product tax of 2. Purchasers' price $=20$; TTM $=20-12=8$; producers' price $=12$;
basic price $=12-2=10$.\\
\textbf{Symbolic.} Consumer pays $p$, trader's cost $q$, product tax $t$, product subsidy $s$:
$\T{TTM}=p-q$; producers' price $=q$; basic price $=q-t+s$.
\end{tmpl}

\begin{tmpl}{Lecture 7 Thai school}
\small A fee of 110 Baht a month from each of 150 students, with a 10\% product tax on the
pre-tax fee, so $110 = 100 + 10$. GVO at producers' prices $=150\times110=16{,}500$; at basic
prices $=150\times100=15{,}000$.\\
\textbf{Symbolic.} $n$ students, fee $f$ including tax at rate $\tau$: producers' $=nf$, basic
$=nf/(1+\tau)$, tax $=nf\tau/(1+\tau)$.
\end{tmpl}

%======================================================================
\section{External transactions and trade exclusions}\label{sec:ext}
\covers{T12}{Lecture 6}{---}

Transactions with the RoW are: exports and imports; primary income (CE and investment income);
current transfers, including taxes on income and wealth; capital transfers; acquisitions less
disposals of valuables and non-produced assets; and financial assets and liabilities.

\begin{itemize}
\item \df{Valuation.} Goods are valued \textbf{f.o.b.} (free on board: excluding insurance and
  freight beyond the exporter's border) at the change of ownership. Customs record imports
  \textbf{c.i.f.} (cost, insurance, freight), but national accounts record imports f.o.b.\ too.
\item \df{Coverage} includes non-residents' spending in the economy (tourists count as exports),
  services such as transport, insurance and hotels, service charges for processing goods owned by
  others, and \df{merchanting}. In merchanting, a merchant in $X$ buys goods from $Y$ and sells
  them to $Z$; ownership changes, but the goods never enter $X$.
\end{itemize}

\begin{keybox}{External trade exclusions (starred in the notes)}
\small \textbf{Rule:} no change of economic ownership between a resident and a non-resident means
no export or import. Excluded:
(1) transactions of non-resident units with the RoW;
(2) flows of samples and goods for exhibition;
(3) goods sent for repair or for rent;
(4) goods of returning residents;
(5) goods sent for processing;
(6) factor services (labour, financial resources), which are primary income, not trade;
(7) goods under operational leasing or rent.
\end{keybox}

\begin{trap}{processing}
\small The \emph{goods} sent abroad for processing are excluded, because the owner never changes.
The processing \emph{fee} is included, as an export of services by the processor's country.
\end{trap}

%======================================================================
\section{Factor incomes, IC, GVA and operating surplus}\label{sec:os}
\covers{T1, T15, T17}{Lectures 5, 7}{2025--26 Q2(i)}

A production process uses labour, inputs (IC), capital goods and natural resources. Part of
the capital goods wears out (depreciation), and part of the non-renewable resources is used up
(depletion). The factors of production are owned, directly or indirectly, by households:

\begin{center}\small
\begin{tabularx}{0.92\textwidth}{@{}l X@{}}
\toprule
\textbf{Factor} & \textbf{Factor payment (compensation)} \\
\midrule
Labour & Compensation of employees, CE (``remuneration of employees'', RE, in 2025 SNA) \\
Land and natural resources & Rent \\
Capital (financial) & Interest \\
Entrepreneurship & Profit \\
A mix of labour and other factors & Mixed income (MI): earnings of unincorporated household enterprises \\
\bottomrule
\end{tabularx}
\end{center}

Rent, interest and profit together make up the \df{operating surplus}.
\begin{itemize}
\item \df{CE / RE}: the total remuneration, in cash or kind, an enterprise owes employees for the
  period's work. It includes wages, salaries, commissions and benefits, \emph{and} the employer's
  social security and PF contributions.
\item \df{Operating surplus (OS)} is the residual once all costs are deducted from output:
  $\T{OS}=\T{GVO}_{bp}-\T{IC}-\T{CE}-\text{other taxes on production (less subsidies)}$. Deduct
  D\&D as well for net OS. OS contains the producer's profit plus the rent and interest paid for
  land and funds owned by others.
\item $\T{GVA}_{bp}$ of corporations and households comes from market output plus output for own
  final use. That of government and NPISHs comes from non-market output plus own final use. For
  the whole economy, $\T{GDP}=\T{GVA}_{bp}+(t-s)_{\text{products}}$.
\end{itemize}

\begin{tmpl}{Lecture 7 bread unit $\to$ asked as 2025--26 Q2(i)}
\small Bread worth 15,000 (without taxes), made from flour 10,000; salaries 1,000; electricity,
fuel and incidentals 2,000.\\
$\T{GVO}_{bp}=15{,}000$; $\T{IC}=10{,}000+2{,}000=12{,}000$; $\T{GVA}_{bp}=3{,}000$;
$\T{OS}=3{,}000-1{,}000=2{,}000$.\\
\textbf{Symbolic} (the exam version: garments $X_1$, fibres $X_2$, salary $X_3$, electricity
etc.\ $X_4$): $\T{GVO}_{bp}=X_1$; $\T{IC}=X_2+X_4$; $\T{GVA}_{bp}=X_1-X_2-X_4$;
$\T{OS}=X_1-X_2-X_3-X_4$.
\end{tmpl}

\begin{trap}{what goes into IC}
\small Salaries are a \emph{factor payment}, never IC. ``Worth (without taxes)'' means basic
prices. \textbf{Factor service, IC or capital good?} Plant and machinery the enterprise owns:
capital good. An office building it hires: IC, since the rental buys a service. Wheat in a flour
mill: IC. Hired land: factor service (rent). A shareholder: factor service. Maintenance of owned
machinery: IC.
\end{trap}

%======================================================================
\section{Depreciation, depletion and gross vs net}\label{sec:dd}
\covers{T10}{Lecture 5}{2025--26 Q1(vi)}

\begin{center}\small
\begin{tabularx}{\textwidth}{@{}l X X@{}}
\toprule
& \textbf{Depreciation} (CFC; called ``depreciation'' in 2025 SNA) & \textbf{Depletion} \\
\midrule
Asset & Produced \emph{fixed} assets & Non-produced \emph{natural} resources (minerals, fuels, wild
stocks) \\
Cause & Wear and tear and normal obsolescence from use in production & Physical removal and
using up by extraction \\
Measures & Cost of fixed capital used up in the period & Natural resources used up in the
period \\
Valued at & \emph{Current} prices of the assets at the time of production, \textbf{not} book
value & Decline in the resource's value \\
In SNA & A cost of production & A cost of production (2025 SNA) \\
\bottomrule
\end{tabularx}
\end{center}

Both reduce the potential to draw benefits from the stock in future. Depreciation can be made good
by new investment; production cannot re-create a depleted natural resource.

\begin{keybox}{Gross $-$ Net $=$ D\&D, for every SNA aggregate}
\centering
NVA $=$ GVA $-$ D\&D; \quad NDP $=$ GDP $-$ D\&D; \quad NNI $=$ GNI $-$ D\&D;\\
net OS $=$ gross OS $-$ D\&D; \quad net profit $=$ gross profit $-$ D\&D; \quad net saving $=$
gross saving $-$ D\&D.
\end{keybox}

\begin{trap}{company depreciation}
\small The depreciation in a company's books is at historic cost, so it is \emph{not} the SNA's
CFC, which is valued at current prices. This is a true/false item in Lecture 12.
\end{trap}
